\documentclass{article}

\usepackage{zharticle}
% \usepackage[screen]{zharticle}  % presentation mode
\input{defs.tex}


\title{\pkg{foils} User Manual}
\author{Hao Zhu}
\date{\today}

\begin{document}
\maketitle

\pkg{foils} is a minimal \pkg{FoilTeX}-inspired LaTeX slide class.
It is implemented as \file{foils.cls}, loads the standard \file{article} class, and adds a small set of commands for writing landscape slide decks with large Computer Modern fonts.

This document describes the public interface provided by the class file \file{foils.cls}.

\section{Overview}
\pkg{foils} is intended for simple lecture or seminar slide decks.
The class sets up letter-sized landscape pages, a large sans-serif default text face, ragged-right paragraphs, slide-friendly list spacing, and a footer that can show the current section and page number.

A typical deck uses the following structure:
\begin{itemize}
  \item \cmd{maketitle} creates the title page.
  \item \cmd{makesection} creates a numbered section page.
  \item \cmd{makefoil} creates each content foil.
  \item \cmd{makeappendix} switches to appendix numbering.
  \item \cmd{makereference} creates a bibliography page.
\end{itemize}

\section{Installation}
For a single talk, copy \file{foils.cls} next to your main \file{.tex} file:
\begin{verbatim}
my-talk/
  talk.tex
  foils.cls
\end{verbatim}

Then select the class in the usual way:
\begin{verbatim}
\documentclass{foils}
\end{verbatim}

For regular use, install \file{foils.cls} in a directory searched by your TeX distribution, such as a local \file{texmf} tree.
After installation, refresh your TeX filename database if your distribution requires it.

\section{Quick start}
The following complete deck \file{talk.tex} uses the main \pkg{foils} commands:
\begin{verbatim}
\documentclass{foils}

\title{A Short Talk}
\author{Firstname Lastname}
\date{\today}
\info{Example Slides}

\begin{document}
\maketitle

\makesection{First section}
\begin{itemize}
  \item First topic
  \item Second topic
\end{itemize}

\makefoil{A first foil}
\begin{itemize}
  \item The Foils class uses large slide-friendly type.
  \item Each foil starts on a new landscape page.
  \item Footers show the current section and page number.
\end{itemize}

\makefoil{A second foil}
\[
  E = mc^2
\]

\makeappendix

\makefoil{Backup material}
Appendix foils are numbered A1, A2, and so on.
\end{document}
\end{verbatim}

Compile the deck \file{talk.tex} with a standard LaTeX engine:
\begin{verbatim}
pdflatex talk.tex
\end{verbatim}

\section{Class options}
\pkg{foils} currently defines one class option:
\begin{description}
  \item[\texttt{lecture}] Reset page numbering at each section and print foil numbers as \texttt{section--page}, for example `\texttt{2--3}'.
  Without this option, normal page numbering is used. Appendix foils use \texttt{A1}, \texttt{A2}, and so on in both modes.
\end{description}
Use the option like this:
\begin{verbatim}
\documentclass[lecture]{foils}
\end{verbatim}

\section{Command reference}
\begin{description}
  \item[\cmd{info}\texttt{\{text\}}] Store metadata shown in small type on section title pages.
    Typical values include a course name, seminar name, or event label.

  \item[\cmd{maketitle}] Create the title page using LaTeX's standard \cmd{title}, \cmd{author}, and \cmd{date} values.
    The title page is numbered internally as page \texttt{0}, so it does not print a normal footer.

  \item[\cmd{makesection}\texttt{\{title\}}] Start a new numbered section on a new page.
    The section title is remembered for later foil footers.
    In \texttt{lecture} mode, the page counter is reset to \texttt{1}.
    The command advances LaTeX's standard \texttt{section} counter, so manual adjustments can use \cmd{setcounter}\texttt{\{section\}\{number\}}.

  \item[\cmd{makefoil}\texttt{\{title\}}] Start a content foil on a new page with a centered title.
    The footer prints the current section title on the left and the page number on the right.

  \item[\cmd{makeappendix}] Start the appendix.
    \pkg{foils} inserts an appendix title page, clears the normal section footer, and numbers later appendix foils as \texttt{A1}, \texttt{A2}, and so on.

  \item[\cmd{makereference}\texttt{[fontsize]\{bibfile\}}] Start a bibliography page using the LaTeX command \cmd{bibliography} and the \texttt{alpha} bibliography style.
    The optional \texttt{fontsize} argument defaults to \cmd{small}; pass another declaration such as \cmd{footnotesize} if the reference list is long.

  \item[\cmd{vf}] Insert stretchable vertical space using \cmd{vspace*}\texttt{\{\cmd{fill}\}}.
    This is useful for pushing material toward the bottom of a foil.
\end{description}

For compatibility with standard LaTeX sectioning commands, \cmd{section} is an alias for \cmd{makesection}, and \cmd{subsection} is an alias for \cmd{makefoil}.
Both aliases accept starred and optional-title forms; the printed title comes from the mandatory argument.

The lower level section commands \cmd{subsubsection}, \cmd{paragraph}, and \cmd{subparagraph} are not supported.
Format deeper in-foil headings manually.

\section{Layout and typography}
\pkg{foils} configures the page as landscape letter paper, with an 11 inch width and an 8.5 inch height.
It sets the PDF page size to the same dimensions.

The default family is sans serif.
Normal body text is set at a large slide-oriented size, with larger display sizes available through the standard
\cmd{large}, \cmd{Large}, \cmd{LARGE}, \cmd{huge}, and \cmd{Huge} declarations.
Paragraphs are ragged right with no indentation, and lists use wider label spacing suitable for projected slides.

The class loads \file{fix-cm}, \file{enumitem}, and \file{graphicx}.
You may load additional packages in your document preamble as usual.

\section{Appendix and references}
Call \cmd{makeappendix} once before backup slides:
\begin{verbatim}
\makeappendix

\makefoil{Additional Details}
This foil is numbered A1.
\end{verbatim}

For bibliographies, create a normal BibTeX file and call \cmd{makereference}:
\begin{verbatim}
\makereference{refs}
\end{verbatim}
This reads entries from \file{refs.bib}.
The class selects the \texttt{alpha} bibliography style.

\section{License}
\pkg{foils} is distributed under the Apache License 2.0.

\end{document}
